\newpage{}
\section{Edge Value Forecasting}
\label{sec:edge-forecast-head}

We use mixture density networks in the sense of Bishop~\cite{bishop1994mdn}, 
to forecast the \texttt{kline} features at each of the edges $e=(u,v) \in \mathcal{E}$%
\note{
  Note, there is separate specialized network for each edge $e=(u,v) \in \mathcal{E}$.
}.  
This means that instead of predicting only one value per feature, we
model a conditional probability density for each future feature $D_{f}$.%
\note{
  This is
  important because the future market state may be multi-modal: several different
  future moves can be plausible under the same current representation.
}

The edge predition heads consumes the edge representation exported by Section~\ref{sec:representation}. 
\[
  \tensor{z}_{u\,v} \in \R^{D_e}.
\]

And tries to estimate the \texttt{klines} at the future: 
\[
  \tensor{f}_{u\,v} \in \R^{C \times H_f \times D_f},
  \quad or \quad
  \tensor{f}_{u\,v}^{c\,i} \in \R^{\times D_f}.
\]

For wich the model creates%
\note{
  A mixture of $K$ Gaussian distributions creates $k = 1,\ldots,K$
  possible centers and spreads for that future, with mixture weights expressing
  how much probability each component carries.
  The density returns one scalar likelihood value for the whole future feature
  vector, not one forecast value per feature.  Feature-level predictions live in
  the component means and scales.  The diagonal covariance assumption factorizes
  each Gaussian component coordinatewise, but the mixture still represents a
  distribution over the full $D_f$-dimensional target vector.
}

\[
  \text{distribution}\quad
  \mathcal{D}_{u\,v}^{c\,i}
  \quad\text{with density}\quad
  {}_{\psi}p_{u\,v}^{c\,i}(\cdot\given\tensor{z}_{u\,v}),
  \quad\text{parameterized by}\quad
  \left\{
    \alpha_{u\,v}^{c\,i\,k},
    \vect{\mu}_{u\,v}^{c\,i\,k},
    \mat{\Sigma}_{u\,v}^{c\,i\,k}
  \right\}_{k=1}^{K}.
\]


Where $\alpha$ is the component weight, $\vect{\mu}$ is the Gaussian mean, and
$\mat{\Sigma}$ the component covariance matrix%
\note{
  The network does not emit every constrained parameter directly ($\alpha$, $\vect{\mu}$, $\mat{\Sigma}$).  
  It emits ($\tilde{\alpha}$,$\vect{\mu}$, $\vect{\beta}$), 
  raw mixture activations $\tilde{\alpha}$, component means $\vect{\mu}$, and raw
  scale vectors $\vect{\beta}$.  The raw activations are normalized into
  weights $\alpha$, and the raw scales are transformed into positive scales
  before forming the covariance matrices $\mat{\Sigma}$.
}.














\newpage{}
\subsection{Edge Forecast Mixture Density Network}
\label{subsec:mdn-forecast}

Let $H_{\psi}\in\N$ be the hidden width of the MDN backbone and let
$R_{\psi}\in\N_{\ge 0}$ be the number of residual Multi-Layer Perceptron blocks.
Each block uses two learned affine maps, layer normalization, and a Sigmoid
Linear Unit activation ($\operatorname{SiLU}$).
The first layer maps one edge representation to hidden space%
\note{
  $\operatorname{SiLU}$ is the coordinatewise Sigmoid Linear Unit activation,
  $\operatorname{SiLU}(x)=x/(1+\exp(-x))$. 
}:
\[
  \vect{h}^{(0)}
  =
  \operatorname{SiLU}
  \left(
    {}_{\psi}\mat{W}_{0}\cdot\tensor{z}_{u\,v}
    +
    {}_{\psi}\vect{b}_{0}
  \right)
  \in \R^{H_{\psi}}.
\]
Each residual block keeps the same width $H_{\psi}$
A typical block has the form
\[
  \vect{q}^{(r)}
  =
  \operatorname{SiLU}
  \left(
    {}_{\psi}\mat{W}^{(r)}_{a_1}\cdot
    {}_{\psi}\operatorname{LN}^{(r)}_{a_1}(\vect{h}^{(r-1)})
    +
    {}_{\psi}\vect{b}^{(r)}_{a_1}
  \right)
  \in \R^{2H_{\psi}},
\]
\[
  \widetilde{\vect{h}}^{(r)}
  =
  {}_{\psi}\mat{W}^{(r)}_{a_2}\cdot\vect{q}^{(r)}
  +
  {}_{\psi}\vect{b}^{(r)}_{a_2}
  \in \R^{H_{\psi}},
\]
and
\[
  \vect{h}^{(r)}
  =
  \operatorname{SiLU}
  \left(
    {}_{\psi}\operatorname{LN}^{(r)}_{a_2}
    \left(
      \vect{h}^{(r-1)}
      +
      \widetilde{\vect{h}}^{(r)}
    \right)
  \right)
  \in \R^{H_{\psi}}.
\]
After the residual stack, the backbone applies layer normalization (${}_{\psi}\operatorname{LN}$)%
\note{
  ${}_{\psi}\operatorname{LN}$ denotes layer normalization applied to the vector inside
  its parentheses.  Superscript $(r)$ labels the residual block, while
  subscripts such as $a_1$ and $a_2$ label different layer normalization
  modules inside that block; they are not tensor
  coordinates.  For a vector $\vect{q}\in\R^{H_{\psi}}$, one common form is
  ${}_{\psi}\operatorname{LN}(\vect{q})
  =
  \vect{g}\odot
  (\vect{q}-\bar{q}\mathbf{1})/\sqrt{s_q^2+\eps}
  +
  \vect{d}$,
  where $\bar{q}=H_{\psi}^{-1}\sum_{a=1}^{H_{\psi}}q^a$,
  $s_q^2=H_{\psi}^{-1}\sum_{a=1}^{H_{\psi}}(q^a-\bar{q})^2$, and
  $\vect{g},\vect{d}\in\R^{H_{\psi}}$ are learned coordinatewise scale and shift
  vectors.
} to the last hidden state as:
\[
  {}_{\psi}\vect{o}
  =
  {}_{\psi}\operatorname{LN}_{\mathrm{out}}
  \left(
    \vect{h}^{(R_{\psi})}
  \right)
  \in \R^{H_{\psi}}.
\]

Based on the shared backbone output ${}_{\psi}\vect{o}$, we apply three
one-layer activationless affine readout branches: one for raw
mixture activations $\tilde{\alpha}$, one for component means $\mu$, and one
for raw scales $\beta$.%
\note{
  The parameters ${}_{\psi}\mat{W}_{(\cdot)}$ and
  ${}_{\psi}\vect{b}_{(\cdot)}$ are the weight and bias of the affine readout
  layer.  The lower-left $\psi$ marks them as MDN
  parameters.
}

After the affine outputs are formed, they are reshaped into channel, horizon,
mixture, and feature axes:

\[
  \begin{aligned}
    \tensor{\tilde{\alpha}}_{u\,v}
    &=
    {}_{\psi}\mat{W}_{\tilde{\alpha}}\cdot{}_{\psi}\vect{o}
    +
    {}_{\psi}\vect{b}_{\tilde{\alpha}}
    \quad \in \R^{C \times H_f \times K},\\
    \tensor{\mu}_{u\,v}
    &=
    {}_{\psi}\mat{W}_{\mu}\cdot{}_{\psi}\vect{o}
    +
    {}_{\psi}\vect{b}_{\mu}
    \quad \in \R^{C \times H_f \times K \times D_f},\\
    \tensor{\beta}_{u\,v}
    &=
    {}_{\psi}\mat{W}_{\beta}\cdot{}_{\psi}\vect{o}
    +
    {}_{\psi}\vect{b}_{\beta}
    \quad \in \R^{C \times H_f \times K \times D_f}.
  \end{aligned}
\]

Equivalently, for every forecast channel $c=1,\ldots,C$, forecast step
$i=1,\ldots,H_f$, and mixture component $k=1,\ldots,K$. For the future features target $\in D_f$ 
these outputs expose:
\[
  \tilde{\alpha}_{u\,v}^{c\,i\,k}\in\R,
  \qquad
  \vect{\mu}_{u\,v}^{c\,i\,k}\in\R^{D_f},
  \qquad
  \vect{\beta}_{u\,v}^{c\,i\,k}\in\R^{D_f}.
\]
The scalar
$\tilde{\alpha}_{u\,v}^{c\,i\,k}$ is the raw activation for mixture component
$k$; the vector $\vect{\mu}_{u\,v}^{c\,i\,k}$ is that component's
$D_f$-dimensional mean; and the vector $\vect{\beta}_{u\,v}^{c\,i\,k}$ is the
raw scale vector that will be transformed into a positive scale.

The collection index $\{c\,i\,k\}$ means that these three quantities are emitted
for every forecast channel, forecast step, and mixture component.  The lower
edge label ${u\,v}$ records that the output was computed from the edge
representation $\tensor{z}_{u\,v}$ for $e=(u,v)$.







\subsection{From Network Outputs to Gaussian Mixtures}
\label{subsec:from-mdn-to-mixtures}

We go from model outputs 
$\{\tensor{\tilde{\alpha}},\,\tensor{\mu},\,\tensor{\beta}\}$
to the distribution's density parameters $\{\alpha,\,\vect{\mu},\,\mat{\Sigma}\}$.\\

For each directed edge $e=(u,v)$, forecast channel $c$, and forecast horizon
$i$, the MDN defines a $K$-component Gaussian mixture over the selected future
target vector $\tensor{f}_{u\,v}^{c\,i}\in\R^{D_f}$.\\

The mixture coefficients%
\note{
  In the diagonal case used here, variances are the squared positive scales.  
  The softmax converts the raw activations $\vect{\tilde{\alpha}}_{u\,v}^{c\,i}$
  into nonnegative mixture coefficients by
  $\softmax(\vect{a})=\exp(\vect{a})/\sum_k\exp(a_k)$.  This guarantees that
  the coefficients form a simplex vector: each component is in $[0,1]$ and the
  components sum to one.
}
are
\[
  \vect{\alpha}_{u\,v}^{c\,i}
  =
  \softmax(\vect{\tilde{\alpha}}_{u\,v}^{c\,i})
  \in [0,1]^K,
  \qquad
  \sum_{k=1}^{K}\alpha_{u\,v}^{c\,i\,k}=1.
\]

The component means are slices of the reshaped mean-head output:
\[
  \vect{\mu}_{u\,v}^{c\,i\,k}
  =
  \left(
    \tensor{\mu}_{u\,v}^{c}
  \right)^{i\,k}
  \in\R^{D_f}.
\]


The raw scale vector is transformed coordinatewise into positive component
scales.  These scales define a diagonal component covariance%
\note{
  The current head emits only coordinatewise scales, so each Gaussian component
  has a diagonal covariance matrix.  This means cross-feature covariance is not
  modeled inside a single component.  The mixture as a whole may still express
  cross-feature structure through different component means and mixture
  coefficients.  A full-covariance MDN would require emitting a
  positive-definite covariance matrix for each component, commonly by predicting
  a Cholesky factor.  That richer parameterization is left as a future
  extension.
}:
\[
  \vect{\sigma}_{u\,v}^{c\,i\,k}
  =
  \softplus(\vect{\beta}_{u\,v}^{c\,i\,k})+\eps\mathbf{1}
  \in\Rpos^{D_f},
  \qquad
  \eps\in\Rpos,
\]
\[
  \mat{\Sigma}_{u\,v}^{c\,i\,k}
  =
  \diag
  \left(
    (\vect{\sigma}_{u\,v}^{c\,i\,k})^{2}
  \right)
  \in\R^{D_f\times D_f}.
\]

Where folows that the predictive density for one future vector is
\[
  {}_{\psi}p_{u\,v}^{c\,i}(\tensor{f}_{u\,v}^{c\,i}\given\tensor{z}_{u\,v})
  :=
  \sum_{k=1}^{K}
  \alpha_{u\,v}^{c\,i\,k}\cdot
  \mathcal{N}
  \left(
    \tensor{f}_{u\,v}^{c\,i};
    \vect{\mu}_{u\,v}^{c\,i\,k},
    \mat{\Sigma}_{u\,v}^{c\,i\,k}
  \right).
\]

So finally, for every edge $e=(u,v)$, forecast channel $c$, and horizon $i$, the MDN
defines one Gaussian predictive distribution with density
${}_{\psi}p_{u\,v}^{c\,i}(\cdot\given\tensor{z}_{u\,v})$%
\note{
  For $\vect{f},\vect{\mu}\in\R^{D_f}$ and positive-definite
  $\mat{\Sigma}\in\R^{D_f\times D_f}$, the Gaussian density is
  $
    \mathcal{N}(\vect{f};\vect{\mu},\mat{\Sigma})
    =
    \left((2\pi)^{D_f}\det(\mat{\Sigma})\right)^{-1/2}
    \exp\!\left(
      -\frac{1}{2}
      (\vect{f}-\vect{\mu})^{\top}
      (\mat{\Sigma})^{-1}
      (\vect{f}-\vect{\mu})
    \right)
  $.
}.

\[
  \mathcal{D}_{u\,v}^{c\,i}
  \in
  \mathcal{P}(\R^{D_f}).
\]





\newpage{}
\subsection{Mixture Density Network's Training Loss}
\label{subsec:mdn-training-loss}

In the training equations below, the leading index $b$ denotes the stacked
edge-labeled sample in the batch%
\note{
  Each batch sample still carries an edge label $(u,v)$.  The edge label is
  suppressed in the loss equations only to keep the batch likelihood readable.
}.
In the coordinatewise likelihood below, $\mu^{b\,c\,i\,k\,d}$ and
$\sigma^{b\,c\,i\,k\,d}$ denote the $d$-th coordinates of the component mean
and positive scale vectors.  For one valid target position $(b,c,i)$, the
unweighted MDN negative log-likelihood is
\[
  \ell^{b\,c\,i}
  =
  -
  \log
  \left[
    \sum_{k=1}^{K}
    \alpha^{b\,c\,i\,k}\cdot
    \prod_{d=1}^{D_f}
    \mathcal{N}
    \left(
      f^{b\,c\,i\,d};
      \mu^{b\,c\,i\,k\,d},
      (\sigma^{b\,c\,i\,k\,d})^{2}
    \right)
  \right].
\]
Equivalently, in log-sum-exp form%
\note{
  For values $a_1,\ldots,a_K$,
  $\operatorname{logsumexp}_{k}(a_k)
  =
  \log\left(\sum_{k=1}^{K}\exp(a_k)\right)$.
  It is used here to combine the $K$ mixture components in log space, which is
  numerically more stable than summing tiny densities directly.
},
\[
  \ell^{b\,c\,i}
  =
  -
  \operatorname{logsumexp}_{k}
  \left(
    \log \alpha^{b\,c\,i\,k}
    +
    \sum_{d=1}^{D_f}
    \log
    \mathcal{N}
    \left(
      f^{b\,c\,i\,d};
      \mu^{b\,c\,i\,k\,d},
      (\sigma^{b\,c\,i\,k\,d})^{2}
    \right)
  \right).
\]
This is the standard maximum-likelihood training objective for the conditional
mixture model.

The implementation may also use feature weights
$\lambda^{(f)}_d\in\Rnonneg$.  In that case the per-component log-density is replaced
by the weighted pseudo-likelihood score
\[
  q^{b\,c\,i\,k}
  =
  \sum_{d=1}^{D_f}
  \lambda^{(f)}_d\cdot
  \log
  \mathcal{N}
  \left(
    f^{b\,c\,i\,d};
    \mu^{b\,c\,i\,k\,d},
    (\sigma^{b\,c\,i\,k\,d})^{2}
  \right),
\]
and
\[
  \ell_{(\lambda)}^{b\,c\,i}
  =
  -
  \operatorname{logsumexp}_{k}
  \left(
    \log \alpha^{b\,c\,i\,k}
    +
    q^{b\,c\,i\,k}
  \right).
\]
When all feature weights are equal to one,
$\ell_{(\lambda)}^{b\,c\,i}$ reduces to the ordinary MDN negative log-likelihood.

Let $\lambda^{(c)}_c\in\Rnonneg$ be an optional channel weight and let
$\lambda^{(h)}_i\in\Rnonneg$ be an optional horizon weight.  The masked batch loss is
\[
  \mathcal{L}_{\mathrm{MDN}}
  =
  \frac{
    \sum_{b=1}^{B}
    \sum_{c=1}^{C}
    \sum_{i=1}^{H_f}
    {}_{f}m^{b\,c\,i}\cdot
    \lambda^{(c)}_c\cdot
    \lambda^{(h)}_i\cdot
    \ell_{(\lambda)}^{b\,c\,i}
  }{
    \sum_{b=1}^{B}
    \sum_{c=1}^{C}
    \sum_{i=1}^{H_f}
    {}_{f}m^{b\,c\,i}\cdot
    \lambda^{(c)}_c\cdot
    \lambda^{(h)}_i
  }.
\]
Only valid future positions contribute to the numerator.  Implementations may
clamp the denominator away from zero for numerical safety.

The scale used inside the loss may also be protected by floors and ceilings,
for example, with $\sigma_{\min}\in\Rpos$,
\[
  \bar{\sigma}^{b\,c\,i\,k\,d}
  =
  \max
  \left(
    \sigma_{\min},
    \sigma^{b\,c\,i\,k\,d}
  \right),
\]
with an optional upper clamp
$\bar{\sigma}^{b\,c\,i\,k\,d}\le\sigma_{\max}$ for
$\sigma_{\max}\in\Rpos$.  These protections do not change the mathematical
object being modeled; they make numerical training stable.





\newpage{}
\subsection{Expectation and covariance observers}
\label{subsec:mdn-expectation-and-covariance}

The expected future vector is available in closed form:
\[
  \widehat{\tensor{f}}_{u\,v}^{c\,i}
  =
  \E_{\psi}
  \left[
    \tensor{f}_{u\,v}^{c\,i}\given\tensor{z}_{u\,v}
  \right]
  =
  \sum_{k=1}^{K}
  \alpha_{u\,v}^{c\,i\,k}\cdot
  \vect{\mu}_{u\,v}^{c\,i\,k}
  \in\R^{D_f}.
\]
In batch form,
\[
  \widehat{\tensor{f}}
  \in
  \R^{B\times C\times H_f\times D_f}.
\]

The full covariance of the mixture distribution is also closed form:
\[
  \operatorname{Cov}_{\psi}
  \left[
    \tensor{f}_{u\,v}^{c\,i}
    \given
    \tensor{z}_{u\,v}
  \right]
  =
  \sum_{k=1}^{K}
  \alpha_{u\,v}^{c\,i\,k}\cdot
  \left(
    \mat{\Sigma}_{u\,v}^{c\,i\,k}
    +
    \vect{\mu}_{u\,v}^{c\,i\,k}(\vect{\mu}_{u\,v}^{c\,i\,k})^{\top}
  \right)
  -
  \widehat{\tensor{f}}_{u\,v}^{c\,i}
  (\widehat{\tensor{f}}_{u\,v}^{c\,i})^{\top}.
\]
Although each component covariance $\mat{\Sigma}_{u\,v}^{c\,i\,k}$ is diagonal, this
mixture covariance can contain off-diagonal terms through the component means.






\subsection{Scalar edge-return marginal}
\label{subsec:mdn-scalar-edge-return-marginal}

Some downstream graph operations use one scalar return coordinate rather than
the full future vector.  Let $d_r\in\{1,\ldots,D_f\}$ be the selected return
coordinate and define
\[
  r_{u\,v}^{c\,i}
  =
  f_{u\,v}^{c\,i\,d_r}
  \in\R.
\]
The scalar predictive density is the marginal mixture
\[
  {}_{\psi}p_{u\,v}^{(r)\,c\,i}(r\given\tensor{z}_{u\,v})
  =
  \sum_{k=1}^{K}
  \alpha_{u\,v}^{c\,i\,k}\cdot
  \mathcal{N}
  \left(
    r;
    \mu_{u\,v}^{c\,i\,k\,d_r},
    (\sigma_{u\,v}^{c\,i\,k\,d_r})^{2}
  \right)
\]
and the corresponding scalar predictive distribution is
\[
  \mathcal{D}_{u\,v}^{(r)\,c\,i}
  \in\mathcal{P}(\R),
  \qquad
  \text{with density }
  {}_{\psi}p_{u\,v}^{(r)\,c\,i}(\cdot\given\tensor{z}_{u\,v}).
\]
Its expectation is
\[
  \widehat{r}_{u\,v}^{c\,i}
  =
  \sum_{k=1}^{K}
  \alpha_{u\,v}^{c\,i\,k}\cdot
  \mu_{u\,v}^{c\,i\,k\,d_r},
\]
and its variance is
\[
  \Var_{\psi}
  \left[
    r_{u\,v}^{c\,i}
    \given
    \tensor{z}_{u\,v}
  \right]
  =
  \sum_{k=1}^{K}
  \alpha_{u\,v}^{c\,i\,k}\cdot
  \left(
    (\sigma_{u\,v}^{c\,i\,k\,d_r})^{2}
    +
    (\mu_{u\,v}^{c\,i\,k\,d_r})^{2}
  \right)
  -
  (\widehat{r}_{u\,v}^{c\,i})^{2}.
\]
These scalar distributions are the objects read by the edge observers in
Section~\ref{sec:distribution-field-and-observer-operators} and by the
edge-to-node recovery step in Section~\ref{sec:edges-to-nodes}.






\subsection{Scope of the head}
\label{subsec:mdn-head-scope}

The MDN head is edge-local.  It defines
\[
  \mathcal{D}_{u\,v}^{c\,i}
  \in\mathcal{P}(\R^{D_f}),
  \qquad
  \text{with density }
  {}_{\psi}p_{u\,v}^{c\,i}(\cdot\given\tensor{z}_{u\,v})
\]
for each directed edge, channel, and forecast step.  It does not by itself define a
joint probability distribution over all active edges, all channels, or all horizons.
When the training loss sums the NLL over $(c,i)$, it is using a pointwise
conditional likelihood over channel-horizon targets.  A joint cross-edge or
cross-horizon uncertainty model is a separate design choice, not part of the
current standardized MDN head.
